\chapter*{Primera parte: Genealogía técnica de la inteligencia artificial}
\addcontentsline{toc}{chapter}{Primera parte: Genealogía técnica de la inteligencia artificial}

\section*{Método a utilizar}
\addcontentsline{toc}{section}{Método a utilizar}

Las historias de la inteligencia artificial suelen contarse como una sucesión de hitos, como un artículo fundacional, un taller, una máquina que asombra a la prensa, un libro que la condena, un invierno, un renacimiento o un modelo que vence a un campeón de un juego. Esa forma de narrar la historia tiene la virtud de la claridad, pero considero que arrastra un supuesto que casi nunca se examina, siendo que los objetos técnicos de la historia -el perceptrón, el sistema experto, la red profunda, el Transformer- son productos acabados que se reemplazan unos a otros según eficacia, como modelos sucesivos de un mismo aparato. Esta primera parte parte de un supuesto distinto tomado de Gilbert Simondon: que un objeto técnico solo puede conococerse adecuadamente si se capta el sentido temporal de su evolución, y que por eso el método adecuado no es la clasificación de objetos ya hechos en géneros y especies, sino el estudio de su génesis. En una nota de El modo de existencia de los objetos técnicos, Simondon lo dice de una forma más precisa. Parafraseo la nota: el punto de usar un método genético es evitar la clasificación que opera después de la génesis, porque las relaciones entre objetos técnicos son tanto horizontales como verticales, y por ello la relación entre ellos es transductiva (Simondon, 2017, p. 26). Conviene subrayar desde ya ese último término, porque es el que esta monografía intentará llevar, en su tercera parte, hasta la idea de una explicabilidad transductiva.
